% MẪU 3 — Chuyên đề ôn thi tốt nghiệp THPT: nhiều chuyên đề, mỗi chuyên đề vài bài.
% Mỗi bài: A. Lý thuyết, B. Ví dụ, C. Bài tập (trắc nghiệm, đúng sai, trả lời ngắn để lẫn).
% Đổi lớp: chỉ sửa dòng \usepackage{toan12}.
\documentclass[12pt]{article}
\usepackage{toan12}
\begin{document}
\loai{3}\chuyende{1}

\tieude{\tentailieu}{Bài 1. Tính đơn điệu của hàm số}

\phan{A. Lý thuyết}
\begin{kienthuc}
Cho hàm số $y=f(x)$ có đạo hàm trên khoảng $K$. Nếu $f'(x)>0$ với mọi $x\in K$ thì $f$ đồng biến trên $K$.
\end{kienthuc}

\phan{B. Ví dụ}
\begin{vd}
Tìm các khoảng đơn điệu của hàm số $y=x^3-3x$.
\giai{$y'=3x^2-3=0\iff x=\pm1$. Hàm số đồng biến trên $(-\infty;-1)$ và $(1;+\infty)$, nghịch biến trên $(-1;1)$.}
\end{vd}

\phan{C. Bài tập}
\begin{tn}[NB]
Hàm số $y=x^3$ đồng biến trên khoảng nào?
\pa{\dung{$\mathbb{R}$}}{$(0;+\infty)$ chỉ}{$(-\infty;0)$}{$(-1;1)$}
\end{tn}
\begin{ds}[TH]
Cho hàm số $y=x^4-2x^2$.
\ydung{$y'=4x^3-4x$.}
\ysai{Hàm số đồng biến trên $(-1;0)$ và $(0;1)$.}
\ydung{Hàm số nghịch biến trên $(-\infty;-1)$.}
\ysai{Hàm số có đúng một điểm cực trị.}
\end{ds}
\begin{tln}[VD]
Có bao nhiêu giá trị nguyên của $m\in[-5;5]$ để hàm số $y=x^3-3x^2+mx$ đồng biến trên $\mathbb{R}$?
\traloi{3}
\loigiai{$y'=3x^2-6x+m\ge0\ \forall x \iff \Delta'=9-3m\le0\iff m\ge3$. Có $m\in\{3;4;5\}$.}
\end{tln}
\begin{tn}[TH]
Hàm số $y=\dfrac{x-1}{x+1}$
\paa{nghịch biến trên $\mathbb{R}\setminus\{-1\}$.}{\dung{đồng biến trên từng khoảng xác định.}}{đồng biến trên $\mathbb{R}$.}{nghịch biến trên từng khoảng xác định.}
\end{tn}

\bangdapan

% ===================== BÀI 2 =====================
\newpage\baimoi
\tieude{}{Bài 2. Cực trị của hàm số}
\phan{A. Lý thuyết}
\phan{B. Ví dụ}
\phan{C. Bài tập}
\end{document}
